% capitulos/capitulado-02-exemplo.tex
%
% Capítulo de conteúdo do MODELO POR CAPÍTULOS AUTÔNOMOS.
% Cada capítulo deste modelo repete a mesma estrutura interna:
% Resumo do Capítulo, Introdução, Objetivo, Hipótese (quando
% aplicável), Metodologia e Contribuições Esperadas. Copie este
% arquivo para criar novos capítulos de conteúdo.
%
% Todo o conteúdo abaixo é fictício e serve apenas para demonstrar o
% padrão estrutural. Substitua pelo conteúdo real da sua pesquisa.

\section{Resumo do Capítulo}
\label{sec:cap2-resumo}

Este capítulo [descreva em duas ou três frases o que o capítulo faz,
qual OE ele endereça e qual hipótese testa, remetendo à
\autoref{sec:oe-h} do \autoref{cap:fundamentacao}].

\section{Introdução}
\label{sec:cap2-introducao}

Ut enim ad minim veniam, quis nostrud exercitation ullamco laboris nisi
ut aliquip ex ea commodo consequat \cite{exemplo:artigo:2024}. Duis aute
irure dolor in reprehenderit in voluptate velit esse cillum dolore eu
fugiat nulla pariatur \cite{exemplo:livro:2020}.

A relevância científica deste capítulo reside na lacuna identificada
na literatura quanto a [descreva a lacuna específica que este capítulo
preenche] \cite{exemplo:capitulo:2019}. Estudos anteriores
\cite{exemplo:anais:2023} apontam que [contextualize o problema
específico deste capítulo], o que motiva o desenho experimental
adotado na \autoref{sec:cap2-metodologia}.

\section{Objetivo}
\label{sec:cap2-objetivo}

\textbf{OE1} (texto completo na \autoref{sec:oe-h}) — [Repita aqui,
em uma frase, o Objetivo Específico 1 já enunciado no
\autoref{cap:fundamentacao}].

\section{Hipótese}
\label{sec:cap2-hipotese}

\textbf{H1} (texto completo na \autoref{sec:oe-h}) — [Repita aqui, em
uma frase, a Hipótese 1 já enunciada no \autoref{cap:fundamentacao}].

\section{Metodologia}
\label{sec:cap2-metodologia}

\subsection{Desenho Amostral}

As coletas foram realizadas em [descreva o número de campanhas, a
sazonalidade e os critérios de amostragem]. Em cada campanha foram
amostrados $n$ pontos distribuídos conforme a
\autoref{fig:pontos-coleta-cap2}.

\begin{figure}[H]
  \centering
  \fbox{\rule{0pt}{6cm}\rule{12cm}{0pt}}
  \caption{Pontos de coleta de amostras na área de estudo. Fonte:
           elaborado pelo autor (2025).}
  \label{fig:pontos-coleta-cap2}
\end{figure}

\subsection{Análises Laboratoriais}

[Descreva aqui os métodos analíticos, equipamentos e protocolos
utilizados, seguindo o mesmo nível de detalhe recomendado pelas
diretrizes do programa] \cite{exemplo:relatorio:2021}.

\subsection{Análise Estatística}

Os dados foram submetidos ao teste de normalidade de Shapiro-Wilk
($\alpha = 0{,}05$). Para comparação entre grupos, utilizou-se
[descreva o teste estatístico adotado], seguindo a abordagem de
\citeonline{exemplo:tese:2023}.

\section{Contribuições Esperadas}
\label{sec:cap2-contribuicoes}

[Descreva aqui os resultados e contribuições esperados deste
capítulo, explicitando como eles alimentam os capítulos
subsequentes e/ou subsidiam a conclusão integrada do trabalho.]
